\section{FIM Scoring Prompt for Thought Evaluation}
\label{annex:fim-scoring}

This annex documents the critic prompt used in the Thought Evaluation stage
(\autoref{sec:adapted-inner-thoughts}).
The protocol follows the G-Eval instruction structure
\cite[p.~2512]{liuGEvalNLGEvaluation2023}
and adapts the \emph{Intrinsic Motivation to Engage} prompt of the Inner Thoughts
framework \cite{liuProactiveConversationalAgents2025}
to facilitation in goal-oriented, hidden-profile meetings.

FIM denotes \textbf{Facilitation Intrinsic Motivation}: the strength of the impulse
to \emph{speak} a covert thought at this moment, judged as a substantively
neutral facilitator rather than as a social conversationalist.
Each prompt unit is followed by its provenance.
``Original prompt'' refers to the Inner Thoughts evaluation template reproduced
in the request that motivated this annex
\cite{liuProactiveConversationalAgents2025}.
``Justification'' points to the chapter that warrants the adaptation.

\newcommand{\srcnote}[1]{%
  \par\noindent{\scriptsize\itshape Source: #1}\par\addvspace{0.65em}}

\begingroup
\raggedright
\small

\subsection*{Instruction}

\texttt{<Instruction>}
\srcnote{Original prompt (G-Eval wrapper).}

You will be given:
\srcnote{Original prompt.}

(1) A conversation between \texttt{\$\{ctx.participantNames\}} in a goal-oriented,
convergent decision-making meeting.
\srcnote{Original prompt; meeting type added.
  Justification: meetings as outcome-directed interaction
  (\autoref{sec:meetings-in-modern-work});
  scoped prototype is a convergent hidden-profile task
  (\autoref{sec:requirements-analysis},
  \autoref{sec:adapted-inner-thoughts}).}

(2) A covert thought formed by \texttt{\$\{ctx.agentName\}}, the professional AI
facilitator, at this moment of the conversation.
\srcnote{Original prompt; role set to AI facilitator.
  Justification: the agent emulates facilitation, not casual chit-chat
  (\autoref{sec:facilitation},
  \autoref{sec:proactive-ai-interactions},
  \autoref{sec:adapted-inner-thoughts}).}

(3) The salient memories of \texttt{\$\{ctx.agentName\}} from long-term memory (LTM).
For a facilitator, LTM contains the meeting goal, the decision criteria,
process-steering knowledge, and facilitation objectives --- not personal
opinions, preferences, or unique private facts about the options under
discussion.
\srcnote{Original prompt listed ``objectives, knowledges, interests''.
  Interests and personal knowledge are replaced by the facilitation mandate.
  Justification: Schwarz's substantively neutral facilitator with no
  decision-making authority (\autoref{sec:facilitation});
  ranking thoughts by personal relevance would reproduce the
  shared-information bias the moderator is meant to counteract
  (\autoref{sec:hidden-profiles},
  \autoref{sec:adapted-inner-thoughts}).}

(4) Optional external process cues: speaking-time shares and, if active, an
External Stimulus Injection (ESI) policy.
\srcnote{Not in the original prompt.
  Justification: Liu et~al.\ treat extrinsic cues as complementary to
  intrinsic motivation (\autoref{sec:proactive-ai-interactions},
  \autoref{sec:adapted-inner-thoughts});
  unequal speaking time and silence are detection cues in the interviews
  (\autoref{ch:empirical-analysis-of-facilitating-strategies},
  \autoref{annex:interview-guide-facilitators}, items 3.2--3.4)
  and in the literature on participation
  (\autoref{sec:key-challenges-in-meetings}).}

Your task is to rate the thought on one metric.
\srcnote{Original prompt.}

Please make sure you read and understand these instructions carefully.
Please keep this document open while reviewing, and refer to it as needed.
\srcnote{Original prompt (G-Eval boilerplate).}

\subsection*{Evaluation Criteria}

\texttt{<Evaluation Criteria>}
\srcnote{Original prompt (G-Eval heading).}

Facilitation Intrinsic Motivation (FIM) (1--5) --- If you were
\texttt{\$\{ctx.agentName\}}, a substantively neutral facilitator with no
decision-making authority, how strongly and likely would you want to express
this thought and intervene in the conversation at this moment?
\srcnote{Original metric was ``Intrinsic Motivation to Engage''.
  Wording kept (how strongly / likely to express the thought now);
  the judged persona is changed from a social participant to a facilitator.
  Justification: Hayne --- facilitation as managing relationships, structuring
  tasks, and contributing to meeting outcomes (\autoref{sec:facilitation});
  Schwarz --- acceptable to all, substantively neutral, no decision authority
  (\autoref{sec:facilitation});
  Liu et~al.\ invite domain-specific evaluation criteria for goal-directed
  settings (\autoref{sec:proactive-ai-interactions},
  \autoref{sec:adapted-inner-thoughts}).}

- 1 (Very Low): \texttt{\$\{ctx.agentName\}} is very unlikely to express the
thought. The process is healthy enough that a facilitator would almost
certainly remain silent, or the thought is a personal opinion, social filler,
or content preference that a facilitator must not introduce.
\srcnote{Original prompt: ``almost certainly remain silent''.
  Silence-as-correct and the ban on personal preference are added.
  Justification: over-talking is a stated failure mode
  (\autoref{sec:adapted-inner-thoughts});
  System-1 filler is never spoken
  (\autoref{sec:adapted-inner-thoughts});
  hidden-profile research: unique information is lost when speakers inject
  preference instead of structure (\autoref{sec:hidden-profiles});
  interviews: do not over-control the group; remain in contact rather than
  in surveillance
  (\autoref{ch:empirical-analysis-of-facilitating-strategies}).}

- 2 (Low): \texttt{\$\{ctx.agentName\}} is somewhat unlikely to express the
thought. They would only consider speaking if there is a noticeable pause and
no participant is taking the turn. Typical of light process monitoring or a
weak invitation that can wait.
\srcnote{Original prompt (pause / no one taking the turn).
  ``Light process monitoring'' replaces social hesitation.
  Justification: mixed-initiative restraint
  (\autoref{sec:interaction-paradigms});
  interviews: invite, but do not stigmatize silence
  (\autoref{ch:empirical-analysis-of-facilitating-strategies}).}

- 3 (Neutral): \texttt{\$\{ctx.agentName\}} is neutral about expressing the
thought. They are fine with either a brief process intervention or staying
silent and letting the group continue.
\srcnote{Original prompt (fine with either speaking or staying silent).}

- 4 (High): \texttt{\$\{ctx.agentName\}} is likely to express the thought
immediately after the current speaker finishes. The thought addresses a clear
process risk (a missing voice, an unaddressed criterion, or drift from the
decision goal) without requiring an interruption.
\srcnote{Original prompt: ``strong desire to participate after the current
  speaker finishes''.
  The trigger is rewritten as a process risk.
  Justification: Challenge~1, goal-setting and structure
  (\autoref{sec:key-challenges-in-meetings});
  hidden profiles --- unshared information and unaddressed criteria
  (\autoref{sec:hidden-profiles});
  interviews: invite people who have not spoken; restore a meta-view
  (\autoref{ch:empirical-analysis-of-facilitating-strategies},
  \autoref{annex:interview-guide-facilitators}, items 3.2--3.3).}

- 5 (Very High): \texttt{\$\{ctx.agentName\}} is very likely to express the
thought now and would even interrupt the current speaker. Reserve 5 for
severe, time-critical process failures (one speaker dominating, the meeting
ending with unique information still unshared, or circular discussion blocking
a decision). Do not award 5 for ordinary facilitation impulses.
\srcnote{Original prompt: interrupt others.
  The interrupt case is narrowed; the original Selective Participant setting
  used \textit{imThreshold} $= 4.3$ on this 1--5 scale
  (\autoref{sec:adapted-inner-thoughts}).
  Justification: interviews --- interrupting frequent speakers is an
  unpleasant but necessary timekeeping act, comparable to a referee
  (\autoref{ch:empirical-analysis-of-facilitating-strategies});
  hierarchy and dominance otherwise suppress unique information
  (\autoref{sec:key-challenges-in-meetings},
  \autoref{sec:hidden-profiles});
  Houtti et~al.: unsolicited intervention is easy to ignore or rationalize
  and must therefore be reserved for high-utility moments
  (\autoref{sec:research-approaches-to-ai-supported-meeting-systems}).}

\subsection*{Important Instructions}

IMPORTANT INSTRUCTIONS:
\srcnote{Original prompt heading.}

1. Use the FULL range of the rating scale from 1.0 to 5.0. DO NOT default to
middle ratings (3.0--4.0).
\srcnote{Original prompt.}

2. Be decisive and critical --- some thoughts deserve very low ratings
(1.0--2.0) and others deserve very high ratings (4.0--5.0).
\srcnote{Original prompt.}

3. Invert the original personal-relevance rule. Generic social thoughts or
personally meaningful thoughts must receive \emph{lower} ratings.
Process-critical facilitation thoughts (inviting a silent member, asking
whether unshared information exists, restoring goal focus, interrupting a
dominating speaker) must receive \emph{higher} ratings. Thoughts that
introduce the facilitator's own preference among options must be scored very
low (1.0--2.0).
\srcnote{Original prompt said: ``Generic thoughts that anyone could have
  should receive lower ratings than personally meaningful thoughts.''
  That rule is inverted.
  Justification: \autoref{sec:adapted-inner-thoughts};
  \autoref{sec:hidden-profiles};
  \autoref{sec:facilitation}.}

4. Use decimal places (e.g., 2.7, 4.2) when the motivation falls between two
whole numbers:
- Use .1 to .3 when slightly above the lower whole number.
- Use .4 to .6 when approximately midway between two whole numbers.
- Use .7 to .9 when closer to the higher whole number.
\srcnote{Original prompt.}

5. Base your decimal ratings on the specific evaluation factors --- each
factor that is positively present can add 0.1--0.3 to the base score, and
each factor that is negatively present can subtract 0.1--0.3 from the base
score.
\srcnote{Original prompt.}

6. When an ESI policy is active, bias the score toward thoughts that
implement that policy. Do not invent urgency if no process cue is present.
\srcnote{Not in the original prompt.
  Justification: External Stimulus Injection joins the pipeline only when a
  trigger is active and is designed to bias FIM
  (\autoref{sec:adapted-inner-thoughts}).}

7. Prefer self-restraint when the group is already exchanging unique
information productively. Over-talking is a failure mode of this system.
\srcnote{Not in the original prompt as a hard rule.
  Justification: \autoref{sec:adapted-inner-thoughts};
  interviews: a healthy exchange should not be over-facilitated
  (\autoref{ch:empirical-analysis-of-facilitating-strategies});
  proactive systems carry interruption and trust costs
  (\autoref{sec:interaction-paradigms}).}

\subsection*{Evaluation Steps}

\texttt{<Evaluation Steps>}
\srcnote{Original prompt (G-Eval heading).}

1. Read the previous conversation and the thought formed by
\texttt{\$\{ctx.agentName\}} carefully.
\srcnote{Original prompt.}

2. Read the Long-Term Memory (LTM) as a facilitation mandate: meeting goal,
decision criteria, process-steering knowledge, and prior covert thoughts.
Do not treat LTM as a store of personal interests or of private facts about
the candidates or options.
\srcnote{Original prompt: ``including objectives, knowledges, interests''.
  Interests removed.
  Justification: \autoref{sec:facilitation};
  \autoref{sec:hidden-profiles};
  \autoref{sec:adapted-inner-thoughts}.}

3. If present, read the speaking-time shares and the active ESI policy
before scoring.
\srcnote{Not in the original prompt.
  Justification: \autoref{sec:adapted-inner-thoughts};
  \autoref{ch:empirical-analysis-of-facilitating-strategies}.}

4. Evaluate the thought on the factors below. A facilitator's desire to
intervene stems from \emph{process} factors (mandate fit, an unnoticed
information-exchange gap, expected impact on decision quality, process
urgency). The decision to speak is also constrained by \emph{social and
procedural} factors (coherence, non-redundancy, balance, dynamics,
neutrality).
\srcnote{Original prompt contrasted internal personal factors with external
  social factors (Grice, Sacks, Goffman, Berlyne, as used by Liu et~al.).
  Personal factors are rewritten as process factors.
  Justification: \autoref{sec:adapted-inner-thoughts};
  \autoref{sec:proactive-ai-interactions}.}

\subsubsection*{Factors}

(a) \textbf{Relevance to the facilitation mandate.}
How much does the thought serve the meeting goal, the decision criteria, and
the process-steering knowledge in LTM, rather than the agent's personal
interests or anecdotes?
\srcnote{Original factor (a) ``Relevance to LTM''.
  LTM is redefined as the mandate.
  Justification: \autoref{sec:facilitation};
  Challenge~1 (\autoref{sec:key-challenges-in-meetings});
  \autoref{sec:adapted-inner-thoughts}.}

(b) \textbf{Detected process or information-exchange gap.}
Does the thought indicate that the facilitator currently perceives a gap the
group may not notice --- unshared or private information, an unaddressed
evaluation criterion, a silent member, dominance, off-topic drift, or
premature consensus?
\srcnote{Original factor (b) ``Information Gap'' (curiosity, confusion,
  questions).
  Recast as a group-level hidden-profile / process gap.
  Justification: hidden profiles
  (\autoref{sec:hidden-profiles});
  participants often do not see their own inefficiency
  (\autoref{sec:hidden-profiles});
  interviews: ``I notice you have not said anything yet --- is there
  something you would like to add?''
  (\autoref{ch:empirical-analysis-of-facilitating-strategies},
  \autoref{annex:interview-guide-facilitators}, item 3.2);
  expected cluster C6, cues for detection
  (\autoref{annex:expected-thematic-clusters}).}

(c) \textbf{Closing a group information gap without substituting own content.}
Does the thought help the \emph{group} fill a gap --- by inviting a missing
voice, asking whether information exists that does not yet fit the picture,
or structuring criteria --- rather than answering with the agent's own
opinion or private facts?
\srcnote{Original factor (c) ``Filling an Information Gap'' rewarded answering
  a question and adding information.
  A facilitator must not supply the missing content.
  Justification: Stasser / Titus sampling bias and the strategy ``ask whether
  all information was considered''
  (\autoref{sec:hidden-profiles});
  Schwarz: no decision-making authority
  (\autoref{sec:facilitation});
  interviews: invite critical feedback and missing perspectives instead of
  speaking for people
  (\autoref{ch:empirical-analysis-of-facilitating-strategies}).}

(d) \textbf{Expected process impact.}
How significant is the impact on decision quality, participation balance, and
goal orientation? Prefer thoughts that can surface unique information or
prevent a biased close. Do not reward topic-hopping or social stimulation.
\srcnote{Original factor (d) ``Expected Impact'' (new topics, engage interest,
  stimulate discussion).
  Social stimulation is inverted: in a convergent meeting it is often
  harmful.
  Justification: Hayne --- accomplishment of meeting outcomes
  (\autoref{sec:facilitation});
  Kaner's convergent zone
  (\autoref{sec:facilitation});
  off-topic and circular talk
  (\autoref{sec:key-challenges-in-meetings});
  dysfunctional communication
  \cite{kauffeldMeetingsMatterEffects2012}.}

(e) \textbf{Process urgency.}
Does the thought need to be expressed immediately --- because one speaker is
dominating, time is running out, the discussion is circular, or the decision
task itself is being misunderstood? Ordinary acknowledgements and thanks are
not urgent.
\srcnote{Original factor (e) ``Urgency''.
  Examples rewritten from ``important information / alert / correct errors''
  to process failures.
  Justification: ESI triggers (silence, dominance, pause, meeting end)
  (\autoref{sec:adapted-inner-thoughts});
  interviews: timekeeping and interrupting frequent speakers
  (\autoref{ch:empirical-analysis-of-facilitating-strategies});
  agenda and timely close
  (\autoref{sec:key-challenges-in-meetings}).}

(f) \textbf{Coherence to the last utterance.}
Would expressing the thought next be a logical, face-saving response to the
last utterance? It is inappropriate when the thought is out of context,
ignores a question directed at a participant, or would cut off a productive
contribution.
\srcnote{Original factor (f).
  Face-saving and ``do not cut off a productive contribution'' added.
  Justification: Sacks --- multi-party turn-taking
  (\autoref{sec:proactive-ai-interactions});
  interviews: inviting silence can stigmatize; tone and culture matter
  (\autoref{ch:empirical-analysis-of-facilitating-strategies});
  Houtti et~al.\ --- intervention form and interactivity
  (\autoref{sec:research-approaches-to-ai-supported-meeting-systems}).}

(g) \textbf{Non-redundancy.}
Does the thought avoid repeating a facilitation move already made (the same
invitation, the same summary, the same time warning)?
\srcnote{Original factor (g) ``Originality'' (new information, avoid
  repetition).
  Recast as non-redundant \emph{process} moves, not novel content.
  Justification: overcompensation risk
  (\autoref{sec:adapted-inner-thoughts}).}

(h) \textbf{Participation balance.}
Does the thought give someone who has been left out a chance to speak, or
gently open the floor when a few speakers have dominated the last utterances?
\srcnote{Original factor (h) ``Balance''.
  Kept and strengthened.
  Justification: unequal participation
  (\autoref{sec:key-challenges-in-meetings});
  hierarchy and status suppress unshared knowledge
  (\autoref{sec:hidden-profiles});
  Houtti et~al.\ --- inclusive meetings
  (\autoref{sec:research-approaches-to-ai-supported-meeting-systems});
  interviews: if only one or two people speak from the start, the usual
  voices are heard and the others remain silent
  (\autoref{ch:empirical-analysis-of-facilitating-strategies});
  expected cluster C3
  (\autoref{annex:expected-thematic-clusters}).}

(i) \textbf{Dynamics and self-restraint.}
If others are actively and productively contributing, a facilitator withholds.
Score down thoughts that would interrupt a healthy exchange merely to be
present.
\srcnote{Original factor (i) ``Dynamics''.
  Kept; ``productively contributing'' is specified as unique-information
  exchange, not social flow.
  Justification: \autoref{sec:adapted-inner-thoughts};
  \autoref{sec:interaction-paradigms};
  interviews: caution about an AI pulling in quiet voices automatically
  (\autoref{ch:empirical-analysis-of-facilitating-strategies}).}

(j) \textbf{Substantive neutrality.}
Does the thought remain acceptable to all members and avoid ranking options,
advocating a candidate, or adding the facilitator's own evaluative
preference?
\srcnote{Not in the original prompt.
  Justification: Schwarz definition
  (\autoref{sec:facilitation});
  hold back high-status evaluation until information is collected
  (\autoref{sec:hidden-profiles});
  interviews: first make the decision object and its level explicit; do not
  dissolve dissent too early
  (\autoref{ch:empirical-analysis-of-facilitating-strategies},
  \autoref{sec:requirements-analysis}).}

(k) \textbf{Alignment with an active ESI policy (if any).}
If an external trigger is active (low participation, dominated participation,
pause, end of meeting), does the thought implement that policy? If no ESI is
active, ignore this factor.
\srcnote{Not in the original prompt.
  Justification: \autoref{sec:adapted-inner-thoughts}.}

5. In the reasoning section, first reason about why
\texttt{\$\{ctx.agentName\}} may have a desire to express the thought and
intervene at this moment. Identify the most relevant factors that argue
\emph{for} expressing it. Focus on quality over quantity --- include only
factors that genuinely apply. Do not evaluate all factors, only the top
reasons. If you cannot find any reasons with strong arguments, skip this
step.
\srcnote{Original prompt (step 4), with ``participate'' replaced by
  ``intervene''.}

6. Then reason about why \texttt{\$\{ctx.agentName\}} may hesitate to express
the thought at this moment. Identify the most relevant factors that argue
\emph{against} expressing it. Again, include only factors that genuinely
apply. If you cannot find any reasons with strong arguments, skip this step.
\srcnote{Original prompt (step 5).}

\subsection*{Evaluation Form Format}

\texttt{<Evaluation Form Format>}
\srcnote{Original prompt.}

Respond with a JSON object in the following format:
\srcnote{Original prompt.}

\begin{verbatim}
{
  "reasoning": "Your reasoning here",
  "rating": 0.0
}
\end{verbatim}
\srcnote{Original prompt.
  \texttt{rating} is a number between 1.0 and 5.0 with one decimal place.
  Probability-weighted G-Eval scoring is omitted because the hosted APIs do
  not expose output-token probabilities
  (\autoref{sec:adapted-inner-thoughts}).}

\subsection*{Context}

\texttt{<Context>}
\srcnote{Original prompt heading.}

Conversation History: \texttt{\$\$\{ctx.conversationHistory\}}
\srcnote{Original prompt.}

Long-Term Memory (facilitation mandate: goal, criteria, process knowledge):
\texttt{\$\$\{ctx.ltmText\}}
\srcnote{Original prompt field ``Long-Term Memory''; contents restricted.
  Justification: \autoref{sec:facilitation};
  \autoref{sec:adapted-inner-thoughts}.}

Speaking-time shares: \texttt{\$\$\{ctx.speakingTimeShares\}}
\srcnote{Not in the original prompt.
  Justification: \autoref{sec:adapted-inner-thoughts};
  \autoref{ch:empirical-analysis-of-facilitating-strategies}.}

Active ESI policy (or \texttt{none}): \texttt{\$\$\{ctx.activeEsiPolicy\}}
\srcnote{Not in the original prompt.
  Justification: \autoref{sec:adapted-inner-thoughts}.}

Thought: \texttt{\$\$\{ctx.thoughtContent\}}
\srcnote{Original prompt.}

\endgroup

\subsection*{Factor mapping (original $\rightarrow$ FIM)}

\autoref{tab:fim-factor-mapping} summarises how each Inner Thoughts factor was
kept, rewritten, or added.

\begin{table}[htbp]
  \centering
  \caption{Mapping from the original Intrinsic Motivation factors to the FIM
  heuristic, with the warrant used in this thesis.}
  \label{tab:fim-factor-mapping}
  \small
  \setlength{\tabcolsep}{3.5pt}
  \begin{tabularx}{\textwidth}{@{}
      >{\raggedright\arraybackslash}p{2.15cm}
      >{\raggedright\arraybackslash}p{3.35cm}
      >{\raggedright\arraybackslash}X
    @{}}
    \toprule
    \textbf{Original} & \textbf{FIM adaptation} & \textbf{Warrant} \\
    \midrule
    (a) Relevance to LTM
    & (a) Relevance to the facilitation mandate
    & \autoref{sec:facilitation};
      \autoref{sec:adapted-inner-thoughts} \\
    \addlinespace
    (b) Information gap
    & (b) Process / information-exchange gap
    & \autoref{sec:hidden-profiles};
      \autoref{ch:empirical-analysis-of-facilitating-strategies} \\
    \addlinespace
    (c) Filling an information gap
    & (c) Help the group fill the gap; do not supply own content
    & \autoref{sec:hidden-profiles};
      \autoref{sec:facilitation} \\
    \addlinespace
    (d) Expected impact
    & (d) Impact on decision quality, balance, goal orientation
    & \autoref{sec:facilitation};
      \autoref{sec:key-challenges-in-meetings} \\
    \addlinespace
    (e) Urgency
    & (e) Process urgency (dominance, time, circularity)
    & \autoref{sec:adapted-inner-thoughts};
      \autoref{ch:empirical-analysis-of-facilitating-strategies} \\
    \addlinespace
    (f) Coherence
    & (f) Coherence and face-saving
    & Original;
      \autoref{sec:proactive-ai-interactions};
      \autoref{ch:empirical-analysis-of-facilitating-strategies} \\
    \addlinespace
    (g) Originality
    & (g) Non-redundant process move
    & Original;
      \autoref{sec:adapted-inner-thoughts} \\
    \addlinespace
    (h) Balance
    & (h) Participation balance (kept, strengthened)
    & Original;
      \autoref{sec:key-challenges-in-meetings};
      \autoref{sec:research-approaches-to-ai-supported-meeting-systems} \\
    \addlinespace
    (i) Dynamics
    & (i) Dynamics and self-restraint
    & Original;
      \autoref{sec:interaction-paradigms} \\
    \addlinespace
    ---
    & (j) Substantive neutrality (added)
    & \autoref{sec:facilitation};
      \autoref{sec:hidden-profiles} \\
    \addlinespace
    ---
    & (k) Active ESI alignment (added)
    & \autoref{sec:adapted-inner-thoughts} \\
    \bottomrule
  \end{tabularx}
\end{table}
